% questions.tex —— 2014 年·辽宁卷·文科 · 题目内容片段（无答案）
% 由 tools/generate_projects.py 依来源总账自 2014/2014_数学真题_辽宁文_试题.md 生成，勿手改（改源 MD 后重跑）


\begin{gknotice}
\textbf{来源说明}：本题面依据固定源 PDF 的两页 A3 题面逐题转录，原卷印刷页码为 897—898。第 897 面为选择题 1—12、填空题 13—16 及解答题 17、18 的题头；第 898 面为解答题 19—21 与选考题 22—24。

\textbf{图形说明}：第 6 题长方形与半圆、第 7 题三视图、第 13 题程序框图、第 19 题四面体、第 20 题圆与切线、第 22 题圆的割线与切线，共 6 幅图由项目内 TikZ 矢量源重绘；其余各题原页无图。

\textbf{解析归属}：固定源不附答案或解析；配套答案与解析均由本库按本题面独立复算。
\end{gknotice}

\gkpart{一、选择题}

\begin{question}
\gkti{1}已知全集 $U=\mathbf R$，$A=\{x\mid x\leqslant 0\}$，$B=\{x\mid x\geqslant 1\}$，则集合 $\complement_U(A\cup B)=$\mbox{（　　）}

（A）$\{x\mid x\geqslant 0\}$；（B）$\{x\mid x\leqslant 1\}$；（C）$\{x\mid 0\leqslant x\leqslant 1\}$；（D）$\{x\mid 0<x<1\}$。

\end{question}
\begin{question}
\gkti{2}设复数 $z$ 满足 $(z-2\mathrm i)(2-\mathrm i)=5$，则 $z=$\mbox{（　　）}

（A）$2+3\mathrm i$；（B）$2-3\mathrm i$；（C）$3+2\mathrm i$；（D）$3-2\mathrm i$。

\end{question}
\begin{question}
\gkti{3}已知 $a=2^{-\frac13}$，$b=\log_2\dfrac13$，$c=\log_{\frac12}\dfrac13$，则\mbox{（　　）}

（A）$a>b>c$；（B）$a>c>b$；（C）$c>a>b$；（D）$c>b>a$。

\end{question}
\begin{question}
\gkti{4}已知 $m$，$n$ 表示两条不同直线，$\alpha$ 表示平面，下列说法正确的是\mbox{（　　）}

（A）若 $m\parallel\alpha$，$n\parallel\alpha$，则 $m\parallel n$；（B）若 $m\perp\alpha$，$n\subset\alpha$，则 $m\perp n$；（C）若 $m\perp\alpha$，$m\perp n$，则 $n\parallel\alpha$；（D）若 $m\parallel\alpha$，$m\perp n$，则 $n\perp\alpha$。

\end{question}
\begin{question}
\gkti{5}设 $\boldsymbol a$，$\boldsymbol b$，$\boldsymbol c$ 是非零向量，已知命题 $p$：若 $\boldsymbol a\cdot\boldsymbol b=0$，$\boldsymbol b\cdot\boldsymbol c=0$，则 $\boldsymbol a\cdot\boldsymbol c=0$；命题 $q$：若 $\boldsymbol a\parallel\boldsymbol b$，$\boldsymbol b\parallel\boldsymbol c$，则 $\boldsymbol a\parallel\boldsymbol c$，则下列命题中的真命题是\mbox{（　　）}

（A）$p\vee q$；（B）$p\wedge q$；（C）$(\neg p)\wedge(\neg q)$；（D）$p\vee(\neg q)$。

\Needspace{0.4\textheight}
\end{question}
\begin{question}
\gkti{6}若将一个质点随机投入如图所示的长方形 $ABCD$ 中，其中 $AB=2$，$BC=1$，则质点落在以 $AB$ 为直径的半圆内的概率是\mbox{（　　）}
\nopagebreak[4]
\begin{center}% ============================================================
% f2014_liaoning_arts_q06_semicircle.tex
% 2014 辽宁卷（文科）第 6 题：长方形 ABCD 与以 AB 为直径的半圆
% 来源 PDF：普通高考/2014/2014辽宁文.pdf
% 原页：第 1 页（A3），印面 897，第 6 题题干下方
% 图中 AB=2，BC=1，半圆在 AB 上方且与 CD 相切。
% ============================================================
\begin{tikzpicture}[>=Stealth, line width=0.7pt, font=\small, x=1cm, y=1cm]
  \coordinate (A) at (0,0);
  \coordinate (B) at (2,0);
  \coordinate (C) at (2,1);
  \coordinate (D) at (0,1);

  % 长方形与半圆
  \fill[gray!25] (A) arc (180:0:1) -- cycle;
  \draw (A) -- (B) -- (C) -- (D) -- cycle;
  \draw (A) arc (180:0:1);

  % 顶点标签
  \fill (A) circle (1.1pt) node[below left] {$A$};
  \fill (B) circle (1.1pt) node[below right] {$B$};
  \fill (C) circle (1.1pt) node[above right] {$C$};
  \fill (D) circle (1.1pt) node[above left] {$D$};
\end{tikzpicture}
\end{center}
\nopagebreak[4]
（A）$\dfrac\pi2$；（B）$\dfrac\pi4$；（C）$\dfrac\pi6$；（D）$\dfrac\pi8$。

\Needspace{0.4\textheight}
\end{question}
\begin{question}
\gkti{7}某几何体三视图如图所示，则该几何体的体积为\mbox{（　　）}
\nopagebreak[4]
\begin{center}% ============================================================
% f2014_liaoning_arts_q07_three_views.tex
% 2014 辽宁卷（文科）第 7 题：三视图
% 来源 PDF：普通高考/2014/2014辽宁文.pdf
% 原页：第 1 页（A3），印面 897，第 7 题题干下方
% 主视图、左视图均为 2×2 矩形，中部有竖线；
% 俯视图为 2×2 正方形去掉左上、右下两个半径 1 的四分之一圆。
% ============================================================
\begin{tikzpicture}[>=Stealth, line width=0.7pt, font=\small,
  dim/.style={line width=0.4pt},
  every node/.style={inner sep=1.5pt}]
  % ---------- 主视图 ----------
  \draw (0,3.0) rectangle (2,5.0);
  \draw (1,3.0) -- (1,5.0);
  \node at (1,2.62) {主视图};
  \draw[dim,<->] (0,5.2) -- (1,5.2) node[midway,above,font=\scriptsize] {$1$};
  \draw[dim,<->] (0,5.6) -- (2,5.6) node[midway,above,font=\scriptsize] {$2$};
  \draw[dim,<->] (-0.3,5.0) -- (-0.3,3.0) node[midway,left,font=\scriptsize] {$2$};

  % ---------- 左视图 ----------
  \draw (3.6,3.0) rectangle (5.6,5.0);
  \draw (4.6,3.0) -- (4.6,5.0);
  \node at (4.6,2.62) {左视图};
  \draw[dim,<->] (3.6,5.2) -- (4.6,5.2) node[midway,above,font=\scriptsize] {$1$};
  \draw[dim,<->] (3.6,5.6) -- (5.6,5.6) node[midway,above,font=\scriptsize] {$2$};

  % ---------- 俯视图 ----------
  % 正方形去掉左上、右下两个四分之一圆（半径 1）
  \draw (0,0) -- (1,0)
        arc[start angle=180, end angle=90, radius=1]
        -- (2,2) -- (1,2)
        arc[start angle=0, end angle=-90, radius=1]
        -- cycle;
  \node at (1,-0.95) {俯视图};
  \draw[dim,<->] (0,-0.32) -- (2,-0.32) node[midway,below,font=\scriptsize] {$2$};
  \draw[dim,<->] (2.35,0) -- (2.35,2) node[midway,right,font=\scriptsize] {$2$};
  % 两个半径标注
  \draw[dim,->] (0,2) -- (0.707,1.293) node[midway,above left,font=\scriptsize] {$1$};
  \draw[dim,->] (2,0) -- (1.293,0.707) node[midway,below right,font=\scriptsize] {$1$};
\end{tikzpicture}
\end{center}
\nopagebreak[4]
（A）$8-\dfrac\pi4$；（B）$8-\dfrac\pi2$；（C）$8-\pi$；（D）$8-2\pi$。

\end{question}
\begin{question}
\gkti{8}已知点 $A(-2,3)$ 在抛物线 $C:y^2=2px$ 的准线上，记 $C$ 的焦点为 $F$，则直线 $AF$ 的斜率为\mbox{（　　）}

（A）$-\dfrac43$；（B）$-1$；（C）$-\dfrac34$；（D）$-\dfrac12$。

\end{question}
\begin{question}
\gkti{9}设等差数列 $\{a_n\}$ 的公差为 $d$，若数列 $\{2^{a_1a_n}\}$ 为递减数列，则\mbox{（　　）}

（A）$d<0$；（B）$d>0$；（C）$a_1d>0$；（D）$a_1d<0$。

\end{question}
\begin{question}
\gkti{10}已知 $f(x)$ 为偶函数，当 $x\geqslant 0$ 时，$f(x)=\begin{cases}\cos\pi x,&x\in\left[0,\dfrac12\right],\\2x-1,&x\in\left(\dfrac12,+\infty\right),\end{cases}$ 则不等式 $f(x-1)\leqslant\dfrac12$ 的解集为\mbox{（　　）}

（A）$\left[\dfrac14,\dfrac23\right]\cup\left[\dfrac43,\dfrac74\right]$；（B）$\left[-\dfrac34,-\dfrac13\right]\cup\left[\dfrac14,\dfrac23\right]$；（C）$\left[\dfrac13,\dfrac34\right]\cup\left[\dfrac43,\dfrac74\right]$；（D）$\left[-\dfrac34,-\dfrac13\right]\cup\left[\dfrac13,\dfrac34\right]$。

\end{question}
\begin{question}
\gkti{11}将函数 $y=3\sin\left(2x+\dfrac\pi3\right)$ 的图象向右平移 $\dfrac\pi2$ 个单位长度，所得图象对应的函数\mbox{（　　）}

（A）在区间 $\left[\dfrac\pi{12},\dfrac{7\pi}{12}\right]$ 上单调递减；（B）在区间 $\left[\dfrac\pi{12},\dfrac{7\pi}{12}\right]$ 上单调递增；（C）在区间 $\left[-\dfrac\pi6,\dfrac\pi3\right]$ 上单调递减；（D）在区间 $\left[-\dfrac\pi6,\dfrac\pi3\right]$ 上单调递增。

\end{question}
\begin{question}
\gkti{12}当 $x\in[-2,1]$ 时，不等式 $ax^3-x^2+4x+3\geqslant 0$ 恒成立，则实数 $a$ 的取值范围是\mbox{（　　）}

（A）$[-5,-3]$；（B）$\left[-6,-\dfrac98\right]$；（C）$[-6,-2]$；（D）$[-4,-3]$。

\end{question}
\gkpart{二、填空题}

\Needspace{0.4\textheight}
\begin{question}
\gkti{13}执行如图的程序框图，若输入 $n=3$，则输出 $T=$\gkblank{}。
\nopagebreak[4]
\begin{center}% ============================================================
% f2014_liaoning_arts_q13_flowchart.tex
% 2014 辽宁卷（文科）第 13 题：程序框图
% 来源 PDF：普通高考/2014/2014辽宁文.pdf
% 原页：第 1 页（A3），印面 897，第 13 题题干下方
% 按原图重绘：循环体依次执行 i=i+1、S=S+i、T=T+S，
% 判断框 $i\leqslant n$ 不成立时输出 T。
% ============================================================
\begin{tikzpicture}[>=Stealth, line width=0.7pt, font=\small,
  every node/.style={inner sep=2pt},
  term/.style={draw, rounded corners=3pt, minimum width=2.1cm,
    minimum height=0.6cm, align=center},
  proc/.style={draw, minimum width=2.4cm, minimum height=0.6cm, align=center},
  io/.style={draw, trapezium, trapezium left angle=70, trapezium right angle=110,
    minimum width=2.4cm, minimum height=0.6cm, align=center},
  dec/.style={draw, diamond, aspect=2, minimum width=2.3cm,
    minimum height=0.9cm, align=center}]
  \node[term] (start) at (0,0) {开始};
  \node[io] (input) at (0,-0.92) {输入 $n$};
  \node[proc] (init) at (0,-1.86) {$i=0,\ S=0,\ T=0$};
  \node[dec] (test) at (0,-3.02) {$i\leqslant n$};
  \node[proc] (b1) at (0,-4.12) {$i=i+1$};
  \node[proc] (b2) at (0,-5.02) {$S=S+i$};
  \node[proc] (b3) at (0,-5.92) {$T=T+S$};
  \node[io] (out) at (0,-7.12) {输出 $T$};
  \node[term] (end) at (0,-8.04) {结束};

  \draw[->] (start) -- (input);
  \draw[->] (input) -- (init);
  \draw[->] (init) -- (test);
  \draw[->] (test.south) -- node[left] {是} (b1.north);
  \draw[->] (b1) -- (b2);
  \draw[->] (b2) -- (b3);
  % 循环体回到判断框左侧
  \draw[->] (b3.west) -- ++(-1.95,0) |- (test.west);
  % “否”分支绕到输出框
  \draw[->] (test.east) -- node[above] {否} ++(1.95,0) |- (out.east);
\end{tikzpicture}
\end{center}
\nopagebreak[4]
\end{question}
\begin{question}
\gkti{14}已知 $x$，$y$ 满足条件 $\begin{cases}2x+y-2\geqslant 0,\\x-2y+4\geqslant 0,\\3x-y-3\leqslant 0,\end{cases}$ 则目标函数 $z=3x+4y$ 的最大值为\gkblank{}。

\end{question}
\begin{question}
\gkti{15}已知椭圆 $C:\dfrac{x^2}9+\dfrac{y^2}4=1$，点 $M$ 与 $C$ 的焦点不重合，若 $M$ 关于 $C$ 的焦点的对称点分别为 $A$，$B$，线段 $MN$ 的中点在 $C$ 上，则 $|AN|+|BN|=$\gkblank{}。

\end{question}
\begin{question}
\gkti{16}对于 $c>0$，当非零实数 $a$，$b$ 满足 $4a^2-2ab+b^2-c=0$ 且使 $|2a+b|$ 最大时，$\dfrac1a+\dfrac2b+\dfrac4c$ 的最小值为\gkblank{}。

\end{question}
\gkpart{三、解答题}

\begin{question}
\gkti{17}在 $\triangle ABC$ 中，内角 $A$，$B$，$C$ 的对边分别为 $a$，$b$，$c$，且 $a>c$，已知 $\overrightarrow{BA}\cdot\overrightarrow{BC}=2$，$\cos B=\dfrac13$，$b=3$，求：

（1） $a$ 和 $c$ 的值；

（2） $\cos(B-C)$ 的值。

\end{question}
\begin{question}
\gkti{18}某大学餐饮中心为了了解新生的饮食习惯，在全校一年级学生中进行了抽样调查，调查结果如下表所示：

\begin{center}
\begin{tabular}{llll}
\toprule
 & 喜欢甜品 & 不喜欢甜品 & 合计 \\
\midrule
南方学生 & $60$ & $20$ & $80$ \\
北方学生 & $10$ & $10$ & $20$ \\
合计 & $70$ & $30$ & $100$ \\
\bottomrule
\end{tabular}
\end{center}

（1） 根据表中数据，问是否有 $95\%$ 的把握认为“南方学生和北方学生在选用甜品的饮食习惯方面有差异”；

（2） 已知在被调查的北方学生中有 $5$ 名数学系的学生，其中 $2$ 名喜欢甜品，现在从这 $5$ 名学生中随机抽取 $3$ 人，求至多有 $1$ 人喜欢甜品的概率。

\Needspace{10\baselineskip}
附：$\chi^2=\dfrac{n(n_{11}n_{22}-n_{12}n_{21})^2}{n_{1+}n_{2+}n_{+1}n_{+2}}$，

\begin{center}
\begin{tabular}{llll}
\toprule
$P(\chi^2\geqslant k)$ & $0.100$ & $0.050$ & $0.010$ \\
\midrule
$k$ & $2.706$ & $3.841$ & $6.635$ \\
\bottomrule
\end{tabular}
\end{center}

\Needspace{0.4\textheight}
\end{question}
\begin{question}
\gkti{19}如图，$\triangle ABC$ 和 $\triangle BCD$ 所在平面互相垂直，且 $AB=BC=BD=2$，$\angle ABC=\angle DBC=120^\circ$，$E$，$F$，$G$ 分别为 $AC$，$DC$，$AD$ 的中点。
\nopagebreak[4]
\begin{center}% ============================================================
% f2014_liaoning_arts_q19_tetrahedron.tex
% 2014 辽宁卷（文科）第 19 题：两个平面垂直的四面体
% 来源 PDF：普通高考/2014/2014辽宁文.pdf
% 原页：第 2 页（A3），印面 898，第 19 题题干下方
% 图中 E、F、G 分别为 AC、DC、AD 的中点；AB、BC、BD 为
% 被遮住的棱，画虚线。
% ============================================================
\begin{tikzpicture}[>=Stealth, line width=0.7pt, font=\small, x=1cm, y=1cm,
  every node/.style={inner sep=1.5pt}]
  \coordinate (A) at (2.00,4.00);
  \coordinate (C) at (4.90,2.20);
  \coordinate (D) at (1.35,0.35);
  \coordinate (B) at (2.75,1.95);
  \coordinate (E) at (3.45,3.10);
  \coordinate (F) at (3.125,1.275);
  \coordinate (G) at (1.675,2.175);

  % 可见棱与截面线
  \draw (A) -- (C) -- (D) -- cycle;
  \draw (E) -- (F);
  \draw (G) -- (C);

  % 被遮住的棱
  \draw[dashed] (A) -- (B);
  \draw[dashed] (B) -- (C);
  \draw[dashed] (B) -- (D);

  % 顶点与中点
  \fill (A) circle (1.1pt) node[above] {$A$};
  \fill (C) circle (1.1pt) node[right] {$C$};
  \fill (D) circle (1.1pt) node[below left] {$D$};
  \fill (B) circle (1.1pt) node[below right] {$B$};
  \fill (E) circle (1.1pt) node[above right] {$E$};
  \fill (F) circle (1.1pt) node[below right] {$F$};
  \fill (G) circle (1.1pt) node[left] {$G$};
\end{tikzpicture}
\end{center}
\nopagebreak[4]
（1） 求证：$EF\perp$ 平面 $BCG$；

（2） 求三棱锥 $D-BCG$ 的体积。

附：锥体的体积公式 $V=\dfrac13Sh$，其中 $S$ 为底面面积，$h$ 为高。

\Needspace{0.4\textheight}
\end{question}
\begin{question}
\gkti{20}圆 $x^2+y^2=4$ 的切线与 $x$ 轴正半轴，$y$ 轴正半轴围成一个三角形，当该三角形面积最小时，切点为 $P$（如图）。
\nopagebreak[4]
\begin{center}% ============================================================
% f2014_liaoning_arts_q20_circle_tangent.tex
% 2014 辽宁卷（文科）第 20 题：圆与第一象限内的切线
% 来源 PDF：普通高考/2014/2014辽宁文.pdf
% 原页：第 2 页（A3），印面 898，第 20 题题干下方
% 圆 x^2+y^2=4 的切线与两坐标轴正半轴围成三角形，切点为 P。
% 图中按面积最小的情形取 P(sqrt2,sqrt2)，切线 x+y=2sqrt2。
% ============================================================
\begin{tikzpicture}[>=Stealth, line width=0.7pt, font=\small, x=1cm, y=1cm,
  every node/.style={inner sep=1.5pt}]
  % 坐标轴
  \draw[->] (-3.1,0) -- (3.9,0) node[right] {$x$};
  \draw[->] (0,-3.1) -- (0,3.9) node[above] {$y$};
  \node[below left] at (0,0) {$O$};

  % 圆
  \draw (0,0) circle (2);

  % 切线与切点
  \draw (2.828,0) -- (0,2.828);
  \fill (1.414,1.414) circle (1.2pt) node[above right] {$P$};
\end{tikzpicture}
\end{center}
\nopagebreak[4]
（1） 求点 $P$ 的坐标；

（2） 焦点在 $x$ 轴上的椭圆 $C$ 过点 $P$，且与直线 $l:y=x+\sqrt3$ 交于 $A$，$B$ 两点，若 $\triangle PAB$ 的面积为 $2$，求 $C$ 的标准方程。

\end{question}
\begin{question}
\gkti{21}已知函数 $f(x)=\pi(x-\cos x)-2\sin x-2$，$g(x)=(x-\pi)\sqrt{\dfrac{1-\sin x}{1+\sin x}}+\dfrac{2x}{\pi}-1$。

（1） 证明：存在唯一 $x_0\in\left(0,\dfrac\pi2\right)$，使 $f(x_0)=0$；

（2） 证明：存在唯一 $x_1\in\left(\dfrac\pi2,\pi\right)$，使 $g(x_1)=0$，且对（1）中的 $x_0$，有 $x_0+x_1>\pi$。

\end{question}
\gkpart{四、选考题（22—24 题，三选一）}

\Needspace{0.4\textheight}
\begin{question}
\gkti{22}【A】如图，$EP$ 交圆于 $E$，$C$ 两点，$PD$ 切圆于 $D$，$G$ 为 $CE$ 上一点且 $PG=PD$，连接 $DG$ 并延长交圆于点 $A$，作弦 $AB$ 垂直 $EP$，垂足为 $F$。
\nopagebreak[4]
\begin{center}% ============================================================
% f2014_liaoning_arts_q22_circle.tex
% 2014 辽宁卷（文科）第 22 题：圆的切线与割线
% 来源 PDF：普通高考/2014/2014辽宁文.pdf
% 原页：第 2 页（A3），印面 898，第 22 题题干下方
% EP 交圆于 E、C，PD 切圆于 D，G 在 CE 上且 PG=PD，
% DG 延长交圆于 A，AB 垂直 EP 于 F。
% 取单位圆：E(-0.866,-0.5)、C(0.866,-0.5)、A(0,-1)、B(0,1)、
% D(0.574,0.819)、P(2.456,-0.5)，此时 PD 与圆相切、PG=PD。
% ============================================================
\begin{tikzpicture}[>=Stealth, line width=0.7pt, font=\small,
  x=1.5cm, y=1.5cm, every node/.style={inner sep=1.5pt}]
  \coordinate (E) at (-0.866,-0.5);
  \coordinate (C) at (0.866,-0.5);
  \coordinate (A) at (0,-1);
  \coordinate (B) at (0,1);
  \coordinate (F) at (0,-0.5);
  \coordinate (G) at (0.158,-0.5);
  \coordinate (D) at (0.574,0.819);
  \coordinate (P) at (2.456,-0.5);

  % 圆
  \draw (0,0) circle (1);

  % 割线 EP、直径 AB、弦 BD、切线 PD、弦 ED、直线 DGA、弦 AC
  \draw (E) -- (P);
  \draw (A) -- (B);
  \draw (B) -- (D);
  \draw (D) -- (P);
  \draw (E) -- (D);
  \draw (D) -- (A);
  \draw (A) -- (C);

  % 点与标签
  \fill (E) circle (1.1pt) node[left] {$E$};
  \fill (C) circle (1.1pt) node[below] {$C$};
  \fill (A) circle (1.1pt) node[below] {$A$};
  \fill (B) circle (1.1pt) node[above] {$B$};
  \fill (F) circle (1.0pt) node[above left] {$F$};
  \fill (G) circle (1.0pt);
  \node at (0.24,-0.62) {$G$};
  \fill (D) circle (1.1pt) node[above right] {$D$};
  \fill (P) circle (1.1pt) node[right] {$P$};
\end{tikzpicture}
\end{center}
\nopagebreak[4]
（1） 求证：$AB$ 为圆的直径；

（2） 若 $AC=BD$，求证：$AB=ED$。

\end{question}
\begin{question}
\gkti{23}【B】将圆 $x^2+y^2=1$ 上每一点的横坐标保持不变，纵坐标变为原来的 $2$ 倍，得曲线 $C$。

（1） 写出 $C$ 的参数方程；

（2） 设直线 $l:2x+y-2=0$ 与 $C$ 的交点为 $P_1$，$P_2$，以坐标原点为极点，$x$ 轴正半轴为极轴建立极坐标系，求过线段 $P_1P_2$ 的中点且与 $l$ 垂直的直线的极坐标方程。

\end{question}
\begin{question}
\gkti{24}【C】设函数 $f(x)=2|x-1|+x-1$，$g(x)=16x^2-8x+1$，记 $f(x)\leqslant 1$ 的解集为 $M$，$g(x)\leqslant 4$ 的解集为 $N$。

（1） 求 $M$；

（2） 当 $x\in M\cap N$ 时，证明：$x^2f(x)+x[f(x)]^2\leqslant\dfrac14$。

\end{question}